% !TEX root = ../main.tex
\chapter{Detection and Segmentation}
\label{ch:detection_segmentation}

Object detection and image segmentation are the two largest
``beyond classification'' workloads in computer vision. Both
take an image and produce a structured output --- a list of
bounding boxes with class labels, or a per-pixel class mask ---
and both have driven the field's hardware demands as much as
classification has. \Lucid\ provides a focused but complete
detection-and-segmentation zoo: Faster~R-CNN
\cite{ren_faster_rcnn_2015}, Mask~R-CNN \cite{he_mask_rcnn_2017},
DETR \cite{carion_detr_2020}, YOLOv3 \cite{redmon_yolov3_2018},
SSD \cite{liu_ssd_2016}, RetinaNet \cite{lin_focal_2017}, FPN
\cite{lin_fpn_2017}, U-Net \cite{ronneberger_unet_2015}, and
DeepLabV3 \cite{chen_deeplab_2017}, with a small set of
panoptic variants \cite{kirillov_panoptic_2019}. This chapter
catalogues the families, develops the mathematics of bounding
boxes / matching / non-maximum suppression, sketches each
architecture, gives Apple-silicon performance figures, and
documents the implementation conventions that the model zoo
follows for these workloads.

\section{Detection: Mathematical Preliminaries}
\label{sec:detection_segmentation:preliminaries}

Before describing models, we set up the mathematical vocabulary
that all detection architectures share. The two universal
primitives are the bounding-box representation, with its IoU
similarity measure, and the non-maximum-suppression
post-processing step.

\subsection{Bounding-box representation}
\label{ssec:detection_segmentation:bbox_repr}

A bounding box is a rectangle in image coordinates. Two common
parameterisations:

\begin{itemize}
  \item \textbf{XYXY}: $(x_1, y_1, x_2, y_2)$ --- the top-left
    and bottom-right corners. Direct for cropping and
    visualisation.
  \item \textbf{CXCYWH}: $(c_x, c_y, w, h)$ --- the centre,
    width, and height. Better numerical conditioning for
    regression because the four components have comparable
    scale.
\end{itemize}

\Lucid\ stores boxes in XYXY for its public-facing
\code{BoundingBox} class (matches COCO \cite{lin_coco_2014}
convention) but converts to CXCYWH internally when training a
regressor. The conversion is
\[
  c_x = \tfrac{1}{2}(x_1 + x_2), \quad
  c_y = \tfrac{1}{2}(y_1 + y_2), \quad
  w = x_2 - x_1, \quad
  h = y_2 - y_1.
\]
The inverse is similarly trivial. Numerical-stability rule: when
$w$ or $h$ approaches zero (degenerate boxes, very small ground
truth), \Lucid\ clamps to a minimum of 1 pixel before computing
logarithms (some loss formulas use $\log w$).

\subsection{IoU (Intersection over Union)}
\label{ssec:detection_segmentation:iou}

The standard similarity measure between two boxes $B_1$ and
$B_2$ is the intersection over union:
\[
  \mathrm{IoU}(B_1, B_2) \;=\;
  \frac{|B_1 \cap B_2|}{|B_1 \cup B_2|}
  \;=\;
  \frac{|B_1 \cap B_2|}{|B_1| + |B_2| - |B_1 \cap B_2|}.
\]
$\mathrm{IoU} = 0$ for non-overlapping boxes, $1$ for identical
boxes, and is symmetric. The intersection area in XYXY
coordinates:
\[
  |B_1 \cap B_2| \;=\;
  \max(0, \min(x_2^{(1)}, x_2^{(2)}) - \max(x_1^{(1)}, x_1^{(2)}))
  \cdot
  \max(0, \min(y_2^{(1)}, y_2^{(2)}) - \max(y_1^{(1)}, y_1^{(2)})).
\]
For training, several extensions are used:
\begin{itemize}
  \item \textbf{GIoU} (Generalised IoU): adds a term that
    penalises the gap between $B_1$ and $B_2$ when they do not
    overlap. Differentiable everywhere.
  \item \textbf{DIoU} (Distance IoU): adds the squared
    centre-distance, normalised by the enclosing-box diagonal.
  \item \textbf{CIoU} (Complete IoU): adds an aspect-ratio
    consistency term on top of DIoU.
\end{itemize}
The choice matters less than the regression loss being
differentiable in the no-overlap region; standard IoU has zero
gradient there. \Lucid\ uses GIoU as the default
\code{box\_regression\_loss}; DIoU and CIoU are available via a
config flag.

\subsection{Non-Maximum Suppression}
\label{ssec:detection_segmentation:nms}

Detectors typically produce many overlapping predictions for the
same object. NMS reduces these to one box per object by:

\begin{enumerate}
  \item Sort predictions by score, descending.
  \item Take the highest-scoring box; add to the output set.
  \item Remove all boxes with $\mathrm{IoU} > \tau$ to the kept
    box.
  \item Repeat until no boxes remain.
\end{enumerate}

The threshold $\tau$ trades recall (low $\tau$ keeps more boxes)
vs precision (high $\tau$ rejects more duplicates); typical
values are $\tau \in [0.5, 0.7]$.

Soft NMS is a variant that, instead of deleting overlapping
boxes, decays their scores:
\[
  s'_i \;=\; s_i \cdot \exp\bigl(-\mathrm{IoU}(B_i, B^*)^2 / \sigma\bigr),
\]
where $B^*$ is the highest-scoring box and $\sigma$ is a
smoothness parameter.

\Lucid\ provides both \code{nms} (hard) and
\code{soft\_nms} as \code{lucid.ops.detection.*}; the hard
version is a fused MPSGraph emitter (fast) while soft NMS uses
the composite path. The hard NMS implementation on Apple
silicon, for batch 64 with 1000 candidates per image, runs in
\SI{2.4}{\milli\second}.

\subsection{Anchor boxes}
\label{ssec:detection_segmentation:anchors}

Anchor-based detectors (Faster~R-CNN, SSD, YOLO, RetinaNet) tile
the image with a dense grid of reference boxes (\emph{anchors})
at multiple scales and aspect ratios. The regression network
predicts displacements from each anchor to the nearest ground
truth, and the classification network predicts a class
probability for each.

Concretely, at each spatial location $(i, j)$ in a feature map
of stride $s$, an anchor of width $w_a$ and height $h_a$ is
centred at $(i \cdot s, j \cdot s)$. The standard parameter
\emph{number of anchors per location} is 3--9.

The training target for an anchor box $(c_x^a, c_y^a, w^a, h^a)$
and ground-truth box $(c_x^g, c_y^g, w^g, h^g)$ is:
\[
  t_x = (c_x^g - c_x^a)/w^a, \quad
  t_y = (c_y^g - c_y^a)/h^a, \quad
  t_w = \log(w^g/w^a), \quad
  t_h = \log(h^g/h^a).
\]
The detector regresses these four scalars; this parameterisation
keeps all four targets at $O(1)$ scale regardless of the
ground-truth box size, which improves training stability.

\section{Faster R-CNN}
\label{sec:detection_segmentation:faster_rcnn}

Faster R-CNN \cite{ren_faster_rcnn_2015} is the canonical
two-stage detector. The pipeline:

\begin{enumerate}
  \item \textbf{Backbone}: a CNN (typically ResNet-50) extracts
    a feature map at stride 16 or 32.
  \item \textbf{Region Proposal Network (RPN)}: a small head on
    the backbone feature map proposes $K$ candidate boxes
    ($K = 1000$ at inference, $K = 2000$ at training).
  \item \textbf{ROI Align}: each proposal extracts a fixed-size
    ($7 \times 7$) sub-feature from the backbone.
  \item \textbf{Detection head}: per-ROI classifier + box
    regressor outputs final detections.
\end{enumerate}

The Region Proposal Network is itself a small detector: a
$3 \times 3$ conv on the backbone feature map, followed by two
sibling $1 \times 1$ convs --- one for objectness scores (2
channels per anchor) and one for box deltas (4 channels per
anchor).

\subsection{Multi-task loss}
\label{ssec:detection_segmentation:faster_rcnn_loss}

The training loss combines RPN and detector losses:
\[
  \mathcal{L}_{\text{Faster R-CNN}} \;=\;
  \mathcal{L}_{\text{RPN}} + \mathcal{L}_{\text{det}},
\]
where each of $\mathcal{L}_{\text{RPN}}$ and
$\mathcal{L}_{\text{det}}$ is itself a sum:
\[
  \mathcal{L} \;=\; \mathcal{L}_{\text{cls}} + \lambda \mathcal{L}_{\text{box}}.
\]
The classification loss is cross-entropy; the box-regression
loss is a smooth-$L_1$ (Huber-style) loss on the four box
deltas. The hyperparameter $\lambda$ balances the two; the
common default is $\lambda = 1$ for the detector and a
class-balance-aware $\lambda$ for the RPN.

\subsection{ROI Align}
\label{ssec:detection_segmentation:roi_align}

ROI Align replaced ROI Pooling (in Fast R-CNN) by using bilinear
interpolation instead of nearest-neighbour quantisation:

\begin{lstlisting}[style=lucid,language=LucidPython]
out = lucid.ops.detection.roi_align(
    feature_map,        # (B, C, H, W)
    boxes,              # (N, 5): [batch_idx, x1, y1, x2, y2]
    output_size=(7, 7),
    spatial_scale=1/16,
    sampling_ratio=2,
)
\end{lstlisting}

The implementation is a fused MPSGraph emitter on Apple silicon
(four bilinear samples per output cell, summed). The backward
is straightforward: each sample contributes to four input
locations with their bilinear weights.

\subsection{Variants}
\label{ssec:detection_segmentation:faster_rcnn_variants}

\code{faster\_rcnn\_resnet\_50\_fpn},
\code{faster\_rcnn\_resnet\_101\_fpn}. Both use the FPN
\cite{lin_fpn_2017} feature pyramid (Section
\ref{sec:detection_segmentation:fpn}) for multi-scale detection.

\section{Mask R-CNN}
\label{sec:detection_segmentation:mask_rcnn}

Mask R-CNN \cite{he_mask_rcnn_2017} extends Faster~R-CNN with a
parallel mask-prediction head. For each ROI, in addition to
classifying and refining its box, the head predicts a binary
mask at $(28, 28)$ resolution per class.

The total loss adds a per-pixel binary cross-entropy:
\[
  \mathcal{L}_{\text{Mask R-CNN}} \;=\;
  \mathcal{L}_{\text{Faster R-CNN}}
  + \mathcal{L}_{\text{mask}},
\]
with $\mathcal{L}_{\text{mask}}$ averaged only over the channel
matching the ground-truth class (decoupling mask and class
predictions).

\Lucid\ provides \code{mask\_rcnn\_resnet\_50\_fpn} and
\code{mask\_rcnn\_resnet\_101\_fpn}. Inference produces both
boxes and per-instance masks.

\section{FPN (Feature Pyramid Network)}
\label{sec:detection_segmentation:fpn}

FPN \cite{lin_fpn_2017} addresses the multi-scale problem in
detection: small objects need fine features, large objects need
coarse features. FPN constructs a feature pyramid in two passes:

\begin{itemize}
  \item \textbf{Bottom-up}: the standard CNN forward pass
    produces features at multiple resolutions (typically
    $C_2$--$C_5$ from ResNet, at strides 4, 8, 16, 32).
  \item \textbf{Top-down}: upsample the coarsest feature, add
    via $1 \times 1$ lateral connection to the next-finer
    backbone feature, repeat. The pyramid levels
    $P_2$--$P_5$ all have 256 channels.
\end{itemize}

The detector heads run on each pyramid level. Small anchors are
assigned to fine levels ($P_2$), large anchors to coarse
levels ($P_5$).

\subsection{Implementation in Lucid}
\label{ssec:detection_segmentation:fpn_impl}

The \code{lucid.models.detection.FPN} module is a separate
class that wraps any backbone with an \code{extract\_features}
method returning a list of multi-scale features. It does not
modify the backbone's training dynamics; the FPN's lateral and
top-down convs are trained jointly with the rest of the
detector.

\section{RetinaNet}
\label{sec:detection_segmentation:retinanet}

RetinaNet \cite{lin_focal_2017} is a single-stage detector that
matched two-stage accuracy by addressing the class-imbalance
problem (most anchors are background, dominating the loss).
The key contribution is the \emph{focal loss}:
\[
  \mathrm{FL}(p_t) \;=\; -\alpha_t (1 - p_t)^\gamma \log(p_t),
\]
where $p_t$ is the model's predicted probability for the
ground-truth class, $\alpha_t$ is a class-balancing weight, and
$\gamma$ (typically 2) reduces the loss contribution from easy
examples (those with $p_t \to 1$).

The intuition: cross-entropy loss is $-\log(p_t)$, which for
$p_t = 0.9$ is still 0.105. With $\gamma = 2$, the focal loss
multiplies by $(1 - 0.9)^2 = 0.01$, reducing the contribution to
$0.001$. Easy examples (the abundant background) thus stop
dominating the gradient.

\Lucid\ provides \code{retinanet\_resnet\_50\_fpn},
\code{retinanet\_resnet\_101\_fpn}.

\section{YOLO Family}
\label{sec:detection_segmentation:yolo}

The YOLO family \cite{redmon_yolo_2016, redmon_yolov3_2018} is
optimised for speed: a single feed-forward pass produces all
detections, with no separate proposal stage. The trade-off is
accuracy (older YOLO versions are below two-stage detectors;
newer ones are competitive).

\subsection{YOLO core idea}
\label{ssec:detection_segmentation:yolo_core}

YOLO divides the image into an $S \times S$ grid; each grid
cell predicts $B$ boxes (and per-class probabilities, jointly
with objectness). At inference, the per-cell predictions are
filtered by score threshold and passed through NMS.

The loss combines:
\begin{itemize}
  \item Box regression on $(c_x, c_y, w, h)$ for grid cells
    containing ground truth.
  \item Objectness binary cross-entropy: 1 for cells with
    ground truth, 0 elsewhere.
  \item Per-class cross-entropy (only on the positive cells).
\end{itemize}

\subsection{YOLOv3 specifics}
\label{ssec:detection_segmentation:yolov3}

YOLOv3 \cite{redmon_yolov3_2018} adds multi-scale prediction
(three grids at strides 32, 16, 8) and a Darknet-53 backbone.
\Lucid\ provides:

\begin{itemize}
  \item \code{yolov3}: full Darknet-53 backbone (61M params).
  \item \code{yolov3\_tiny}: paper-cited mobile variant
    (8.7M params). The paper itself names \emph{tiny}, so H11 is
    satisfied.
\end{itemize}

\section{SSD (Single Shot MultiBox Detector)}
\label{sec:detection_segmentation:ssd}

SSD \cite{liu_ssd_2016} predates RetinaNet as a single-stage
multi-scale detector. The backbone (typically VGG-16) is
augmented with extra conv stages at progressively smaller
resolutions; each scale predicts a fixed set of anchor boxes.

\code{ssd\_300}: $300 \times 300$ input, VGG-16 backbone, six
prediction scales.
\code{ssd\_512}: $512 \times 512$ input, six prediction scales
(SSD512, paper variant).

SSD is less commonly used today (YOLOv3+ and RetinaNet
outperform it) but is included for completeness and for older
research baselines.

\section{DETR (DEtection TRansformer)}
\label{sec:detection_segmentation:detr}

DETR \cite{carion_detr_2020} reformulates detection as a
direct set-prediction problem, eliminating both anchors and NMS.
The pipeline:

\begin{enumerate}
  \item CNN backbone (ResNet-50) extracts a feature map.
  \item Transformer encoder processes the flattened features
    with positional embeddings.
  \item Transformer decoder consumes $N$ \emph{object queries}
    (learnable embeddings) and the encoder output; each query
    becomes one detection.
  \item Each decoder output is projected to a class label and
    a bounding box.
\end{enumerate}

The set of $N$ detections (default $N = 100$, much larger than
the typical object count) is matched to ground truth via the
Hungarian algorithm. Unmatched detections are labelled
``no-object''.

\subsection{The Hungarian matching}
\label{ssec:detection_segmentation:hungarian}

Let predictions be $\{(\hat{p}_i, \hat{b}_i)\}_{i=1}^N$ and
ground truths be $\{(c_j, b_j)\}_{j=1}^M$ (with $M \leq N$).
Define the matching cost:
\[
  \mathcal{C}(i, j) \;=\; -\hat{p}_i(c_j) + \lambda_1 \|b_j - \hat{b}_i\|_1 + \lambda_2 \mathcal{L}_{\text{GIoU}}(b_j, \hat{b}_i).
\]
The Hungarian algorithm finds the permutation $\sigma$
minimising $\sum_j \mathcal{C}(\sigma(j), j)$ in $O(N^3)$ time.
For $N = 100$, this is microseconds per image and not a
bottleneck.

\subsection{Training loss}
\label{ssec:detection_segmentation:detr_loss}

Given the matching $\sigma$:
\[
  \mathcal{L}_{\text{DETR}} \;=\;
  \sum_{j=1}^M \bigl[
    -\log \hat{p}_{\sigma(j)}(c_j)
    + \lambda_1 \|b_j - \hat{b}_{\sigma(j)}\|_1
    + \lambda_2 \mathcal{L}_{\text{GIoU}}(b_j, \hat{b}_{\sigma(j)})
  \bigr]
  + \sum_{i \notin \text{im}(\sigma)} -\log \hat{p}_i(\varnothing).
\]
The second term penalises detections matched to the no-object
class. The class loss is weighted lightly to avoid the
no-object class dominating.

\subsection{Variants}
\label{ssec:detection_segmentation:detr_variants}

\code{detr\_resnet\_50}, \code{detr\_resnet\_101}.
\Lucid\ does not yet provide the Deformable DETR or DINO
variants; these are tracked for a future release.

\section{U-Net}
\label{sec:detection_segmentation:unet}

U-Net \cite{ronneberger_unet_2015} is the foundational
encoder--decoder architecture for semantic segmentation. The
shape resembles a U: progressive downsampling (encoder) then
progressive upsampling (decoder), with \emph{skip connections}
that copy each encoder stage to the matching decoder stage.

The skip connections concatenate (not add) the encoder
features to the decoder features along the channel axis, then
two $3 \times 3$ convolutions mix them. The motivation:
upsampling alone loses spatial detail; the skip carries the
fine-grained spatial information that the encoder still has.

\Lucid\ provides \code{unet} (the original biomedical
configuration), \code{unet\_resnet\_50} (ResNet backbone +
U-Net decoder, a common practical variant).

\section{DeepLabV3}
\label{sec:detection_segmentation:deeplab}

DeepLab \cite{chen_deeplab_2017} addresses the same semantic
segmentation problem with a different design: a CNN backbone
with \emph{atrous} (dilated) convolutions that preserve spatial
resolution at the cost of more compute, followed by an
\emph{Atrous Spatial Pyramid Pooling} (ASPP) module that captures
multi-scale context.

The dilated convolution:
\[
  y(i) \;=\; \sum_k x(i + r \cdot k) \cdot w(k),
\]
where $r$ is the dilation rate. With $r = 2$, a $3 \times 3$
kernel covers a $5 \times 5$ receptive field with $3 \times 3$
parameters.

\Lucid\ provides \code{deeplabv3\_resnet\_50},
\code{deeplabv3\_resnet\_101}.

\section{Panoptic Segmentation}
\label{sec:detection_segmentation:panoptic}

Panoptic segmentation \cite{kirillov_panoptic_2019} unifies
instance segmentation (separating object instances) and semantic
segmentation (per-pixel labelling) into a single output. The
metric is panoptic quality (PQ):
\[
  \mathrm{PQ} \;=\;
  \underbrace{\frac{\sum_{(p,g) \in TP} \mathrm{IoU}(p,g)}{|TP|}}_{\text{Segmentation Quality}}
  \cdot
  \underbrace{\frac{|TP|}{|TP| + \tfrac12|FP| + \tfrac12|FN|}}_{\text{Recognition Quality}}.
\]
\Lucid\ provides \code{panoptic\_fpn\_resnet\_50} (a
panoptic-FPN variant that adds a semantic-segmentation head on
top of Mask~R-CNN).

\section{Performance on Apple Silicon}
\label{sec:detection_segmentation:perf}

Per-iteration training time on M3~Max for selected
detection/segmentation models, ImageNet pretraining + COCO
fine-tuning, batch 4, $800 \times 1024$ input, FP32:

\begin{table}[t]
\centering
\small
\begin{tabular}{lrr}
\toprule
Model & Params (M) & Step time (ms) \\
\midrule
\code{faster\_rcnn\_resnet\_50\_fpn}  & 41.6 & 380 \\
\code{mask\_rcnn\_resnet\_50\_fpn}    & 44.2 & 470 \\
\code{retinanet\_resnet\_50\_fpn}     & 38.0 & 295 \\
\code{yolov3}                         & 61.5 & 220 \\
\code{ssd\_300}                       & 24.0 & 130 \\
\code{detr\_resnet\_50}               & 41.0 & 410 \\
\code{unet}                           & 7.8  & 95  \\
\code{deeplabv3\_resnet\_50}          & 42.0 & 310 \\
\bottomrule
\end{tabular}
\caption[Detection/segmentation throughput.]{Per-step training
time on M3~Max for selected detection and segmentation
architectures. YOLO is the fastest single-stage detector; DETR
is competitive but slower due to attention. Mask R-CNN is
slower than Faster R-CNN by the mask-head overhead (about 25\%).
The two-stage detectors include the per-batch overhead of
proposal-set padding and ROI Align.}
\label{tab:detection_segmentation:perf}
\end{table}

The numbers omit data-loading time (handled separately, see
\chref{ch:data_pipeline}). The compile path
(\chref{ch:compile_design}) is currently disabled for
detection models because of dynamic-shape requirements (variable
number of proposals per image); the 3.5.1 release with dynamic
shapes will enable it.

\section{Common Implementation Patterns}
\label{sec:detection_segmentation:patterns}

\subsection{The data format}
\label{ssec:detection_segmentation:data_format}

Detection targets are a list of dicts:

\begin{lstlisting}[style=lucid,language=LucidPython]
target = {
    "boxes": Tensor(shape=(N, 4), dtype=fp32),  # XYXY
    "labels": Tensor(shape=(N,), dtype=int64),
    "masks": Tensor(shape=(N, H, W), dtype=bool),  # optional
}
\end{lstlisting}

The dataset returns one such dict per image. The
\code{DataLoader} (\chref{ch:data_pipeline}) collates a batch
into the format the model expects (typically a padded tensor of
boxes plus a per-image valid count).

\subsection{The backbone interface}
\label{ssec:detection_segmentation:backbone_iface}

Every detection model takes a backbone via a small interface:

\begin{lstlisting}[style=lucid,language=LucidPython]
class DetectionBackbone(nn.Module):
    out_channels: list[int]
    out_strides: list[int]
    def forward(self, x: Tensor) -> list[Tensor]: ...
\end{lstlisting}

The backbone returns features at multiple scales (typically
4 levels). FPN wraps any backbone to standardise the channel
counts at 256 across all levels.

\subsection{The training-vs-inference output}
\label{ssec:detection_segmentation:train_vs_infer}

Detection models have asymmetric outputs:
\begin{itemize}
  \item \textbf{Training}: returns the loss dict (per-component
    losses for logging).
  \item \textbf{Inference}: returns a list of detections per
    image (boxes, labels, scores).
\end{itemize}
The same \code{model(x, targets=...)} call switches based on
whether \code{targets} is passed. This matches torchvision's
convention.

\section{Pretrained Weights}
\label{sec:detection_segmentation:pretrained}

Detection models are pretrained on COCO \cite{lin_coco_2014}.
The standard entry points:

\begin{lstlisting}[style=lucid,language=LucidPython]
model = lucid.models.AutoModel.from_pretrained(
    "faster_rcnn_resnet_50_fpn",
    pretrained_on="coco",
)
# Returns a model ready for COCO inference (91 classes including bg)
\end{lstlisting}

For fine-tuning on a custom dataset, replace the classifier head
with a new one matching the number of classes:

\begin{lstlisting}[style=lucid,language=LucidPython]
num_classes = 10  # Your dataset
in_features = model.roi_heads.box_predictor.cls_score.in_features
model.roi_heads.box_predictor = lucid.models.detection.FastRCNNPredictor(
    in_features, num_classes
)
\end{lstlisting}

\section{Summary and Forward References}
\label{sec:detection_segmentation:summary}

The detection-and-segmentation zoo in \Lucid: Faster~R-CNN,
Mask~R-CNN (two-stage), RetinaNet, YOLOv3, SSD (single-stage),
DETR (set-prediction), U-Net, DeepLabV3 (semantic segmentation),
panoptic-FPN (panoptic). Each is paired with FPN or a similar
multi-scale module when needed; each follows the same backbone
interface so the zoo's components are interchangeable. COCO
pretrained weights are available through the hub.

The next chapter,
\chref{ch:language_models}, treats language models. The
\chref{ch:generative} chapter handles generative models
(GAN/VAE/diffusion). \chref{ch:pretrained} treats the weights
hub itself.

\begin{lucidnote}
For the user: \code{mask\_rcnn\_resnet\_50\_fpn} is the
standard ``works-out-of-the-box'' object-and-instance detector;
for speed-critical workloads, prefer \code{yolov3}. For
semantic segmentation, \code{deeplabv3\_resnet\_50} is
recommended; \code{unet} is a small alternative for biomedical
imaging.
\end{lucidnote}
